\section{Recomendaciones}
\begin{indentedsubsection}
Para garantizar el éxito comercial y logístico a largo plazo de la plataforma Odoo en Industria Textil Atlas E.I.R.L., se sugieren las siguientes recomendaciones operativas y estratégicas:
\end{indentedsubsection}

\subsection{Disciplina y exclusividad en el registro de información}
\begin{indentedsubsection}
Para mantener la utilidad y veracidad de los reportes financieros del negocio, es imperativo que todo el personal de tienda registre las ventas de mostrador en el Punto de Venta (POS) y las cotizaciones mayoristas en Ventas en tiempo real. Se debe desterrar por completo el uso paralelo de libretas físicas o registros informales fuera del sistema. La precisión del inventario físico y la detección automática de reabastecimiento dependen exclusivamente de que no existan operaciones ``ocultas'' fuera de Odoo.
\end{indentedsubsection}

\subsection{Unificación y escalabilidad hacia la Facturación Electrónica (SUNAT)}
\begin{indentedsubsection}
De acuerdo con el alcance técnico inicial, la facturación formal con SUNAT continuará gestionándose a través del sistema facturador externo actual para no recargar la curva de aprendizaje inicial del personal en tienda. Sin embargo, como recomendación clave para el mediano plazo, se aconseja planificar la integración directa del facturador electrónico dentro de Odoo. Esto evitará mantener un doble sistema (Odoo para el registro de ventas/POS e inventarios, y un sistema externo para emitir comprobantes tributarios), eliminando la duplicidad de tareas y la necesidad de transcribir información manualmente. Con la unificación total, las boletas y facturas con validez fiscal se emitirán automáticamente desde la misma orden de venta ganada o el POS, cerrando el ciclo comercial y fiscal en una sola plataforma.
\end{indentedsubsection}

\subsection{Aprovechamiento activo de reglas de reorden}
\begin{indentedsubsection}
Se recomienda a la gerencia auditar y ajustar de forma periódica las reglas de stock mínimo y reabastecimiento automático en base a la estacionalidad de las ventas de artículos deportivos (por ejemplo, incrementando los stocks de seguridad durante campañas escolares o campeonatos locales). Esto evitará la acumulación innecesaria de existencias de baja rotación y garantizará la disponibilidad permanente de productos de alta demanda en mostrador.
\end{indentedsubsection}

\subsection{Resguardo de datos comerciales en la nube}
\begin{indentedsubsection}
Dado que el catálogo de productos con sus costos reales de adquisición de talleres y los precios de venta son activos estratégicos de la empresa, se recomienda aprovechar la infraestructura segura de Odoo en la nube. Se sugiere restringir los permisos de usuario de ventas para evitar fugas de información sensible sobre márgenes y costos de proveedores, y mantener centralizado todo el historial de transacciones en la plataforma para evitar pérdidas de información por fallas en equipos locales o teléfonos móviles de los vendedores.
\end{indentedsubsection}
